\documentclass[11pt]{article}

\usepackage[margin=1in]{geometry}
\usepackage[T1]{fontenc}
\usepackage[utf8]{inputenc}
\usepackage{lmodern}
\usepackage{microtype}
\usepackage{amsmath,amssymb}
\usepackage{xcolor}
\usepackage{enumitem}
\usepackage{array}
\usepackage{tabularx}
\usepackage{booktabs}
\usepackage{hyperref}

\hypersetup{
  colorlinks=true,
  linkcolor=blue!45!black,
  urlcolor=blue!45!black,
  citecolor=blue!45!black,
  pdftitle={Foundational Conjecture},
  pdfauthor={Hoapa project}
}

\setlist[itemize]{leftmargin=1.6em,itemsep=0.35em,topsep=0.35em}
\setlength{\parindent}{1.2em}
\setlength{\parskip}{0pt}
\sloppy

\newcolumntype{Y}{>{\raggedright\arraybackslash}X}

\title{Motivation for the Arc\\[0.35em]\large The Foundational Conjecture}
\author{Hoapa project}
\date{June 2026}

\begin{document}
\maketitle

\section*{The foundational conjecture}

That modalities (perspectivalism encoded in HoTT) allows complexity to develop
is trivial in abstract math [ModalHoTT].

That typed data is more robust, readable, literate\ldots{} is trivial in any
domain concerned with information transfer, such as programming (HoTT says in
homotopic math as well) [SessionValues].

That unityped agents are most adapted to handle complexity is trivial in
concurrency theory\footnote{And it is the basis of the actor model (and only in
this relatively new domain because it is in friction with a foundational
assumption of modernity, that complexity is based on specialisation).} (a
lesson of the Yoneda lemma, visible in the notion of comonads/coalgebras?)
[ActorBase].

The result we're targeting should show (encoded in category theory):

\paragraph{If:}
\begin{itemize}
\item SessionValues and ActorBase follow separation of concern (independence,
in the sense of modularity/clean code).

\item ModalHoTT and SessionValues--ActorBase are in a requires--allows
adjunction.

\emph{codependency:}
\begin{itemize}
\item perspectives require (the optimality) typed data handled by unityped (or
``univocal'' in spinozism) agents.
\item persepctives allow typed data to be handled by unityped agents.
\item the separation of concern between data and agents (more precisely, the
optimality of the system, secured by that separation) requires
perspectivalism (a modal network).
\item it allows it.
\end{itemize}
\end{itemize}

\paragraph{Then:}
Philosophically: the category which hosts this structure is capable of ?
(World understanding and world creation)

Mathematically: something about covers/coverings, recursion, a form of
contribution to a developing infinite network\ldots{}

\par\medskip
The last item is the mathematical fragment, interesting for its own sake in
mathematics, of a foundational statement in the space of interdisciplinary
theories, probably similar to the Yoneda lemma.

\clearpage
\section*{Current first checked fragment}

On the operational side, the project now has a point-generated sheafification
reflector and a sheaf-level Benton/LNL adjunction over the supported
event-discrete actor-causal site. This is not yet the cross-site \(Allow \dashv
Require\) adjunction; it is the first modest B-internal theorem that makes one
side of that future adjunction concrete.

\section*{The cascade of trade-offs: a praxis response to structural problems}

The framework is, read operationally, a sequence of engineering trade-offs that
the foundational commitments generate, together with the strategies that
resolve them. The cascade is \emph{generative}: each resolution opens the next
problem. The strategy applied at every step is the same shape---separation of
concerns. The full layer-by-layer catalogue is in \texttt{future\_work.md};
what follows is the compact picture.

\begin{table}[htbp]
\small
\setlength{\tabcolsep}{5pt}
\renewcommand{\arraystretch}{1.18}
\begin{tabularx}{\linewidth}{@{}>{\bfseries}p{0.22\linewidth}Y Y@{}}
\toprule
Layer & Trade-off it navigates & Strategy \\
\midrule
Hybrid &
Locality (actors operate in frames) vs globality (content must cohere across
frames) &
Sheaf descent: compatible local content glues to global \\
\addlinespace
Reflexive (\(\Pi\)) &
Stability (the protocol must be reliable) vs adaptability (it must respond to
change) &
Meta-conversations propose, consensus ratifies, non-retroactivity protects the
past \\
\addlinespace
Bulk (\(\Lambda\)) &
Expressive power (rich vocabulary) vs transparency (decodable by every
participant) &
Protocol generates its own language; the reflexive layer maintains its
transparency over time \\
\addlinespace
Field-theoretic dynamics &
Exploration (proposal storm, high entropy) vs exploitation (stable mature
language, low entropy) &
RG flow with reflexive insertions; symmetry-breaking \(\to\) fluctuation
\(\to\) stabilization phases \\
\addlinespace
Federation &
Coherence within one protocol vs plurality across protocols &
Typed translation interfaces, partial clones, conflict-as-data \\
\addlinespace
Information cohomology &
Dense encoding (expert-economized) vs accessible encoding (decodable by the
unityped receiver population) &
Receiver-relative measurement of encoding cost; transparency taken as an
engineering -ility on the same footing as performance and security \\
\bottomrule
\end{tabularx}
\end{table}

The cascade is not a list of features added by choice. It is the sequence of
problems the foundational commitments create, with each layer being the
structural response to the problem the previous layer's resolution exposed.
Hybrid's locality-vs-globality only becomes a problem because we committed to
actors and sessions as separate legs. Reflexive's stability-vs-adaptability
only becomes a problem because Hybrid gave us a maintainable artifact. Bulk's
expressivity-vs-transparency only becomes a problem because Reflexive gave us a
self-editing protocol that must now talk about things. Field-theoretic
dynamics' exploration-vs-exploitation only becomes a problem because Bulk gave
us a language that must evolve well. Federation's coherence-vs-plurality only
becomes a problem because the four previous layers gave us a complete protocol
that, in isolation, would re-create the one-true-framework failure the project
was built to dissolve. Information cohomology's dense-vs-accessible encoding
becomes salient as soon as Bulk produces a language, and acute as soon as
Federation puts that language in front of multiple receiver populations.

The strategy is uniform across the cascade. Every trade-off is resolved by
separating the two concerns that are in tension, giving each its own
categorical home, and providing the precise interface between them. Locality
and content are separated by the sheaf; current protocol and its edits are
separated by \(\Pi\); how and what are separated by the session/content split;
content and dynamics are separated by the factorization-algebra-plus-RG-flow;
interior and exterior of a protocol are separated by the meta-topos; content
and its measurement are separated by the (co)homology pairing. The foundational
commitments themselves---actor-vs-content on B, spatial-vs-discrete on A
(cohesion), context-vs-system at the cross-site adjunction---are the same
architectural move at the foundational floor.

So separation of concerns is not an engineering principle the framework
happens to use. It is the framework's primitive architectural move, applied
recursively from the foundations upward, and the cascade of trade-offs is what
makes that recursion necessary. The transparency-of-language trade-off
(information cohomology) belongs on this list on the same footing as the
others---not relegated to ``measurement''---because it is the encoding question
the unityped commitment forces on the engineer: \emph{how is the unityped
receiver going to understand the message?} That question is no less an
engineering -ility than performance, security, or maintainability. It is just
not currently on the engineer's radar, because the economy of expertise prefers
density to transparency. The framework's commitment is to put it there.

\par\bigskip
\noindent\emph{Authorship note. The first section was authored by a human. The following two
sections were authored by AI agents.}

\end{document}
